\documentclass[conference]{IEEEtran}
\IEEEoverridecommandlockouts

\usepackage[T1]{fontenc}
\usepackage[utf8]{inputenc}
\usepackage[english]{babel}
\usepackage{cite}
\usepackage{amsmath,amssymb,amsfonts}
\usepackage{booktabs}
\usepackage{array}
\usepackage{graphicx}
\usepackage{url}
\usepackage[hidelinks]{hyperref}
\newcommand{\code}[1]{\textnormal{#1}}

\setlength{\textfloatsep}{7pt plus 1pt minus 2pt}
\setlength{\floatsep}{5pt plus 1pt minus 2pt}
\setlength{\intextsep}{5pt plus 1pt minus 2pt}
\setlength{\abovedisplayskip}{4pt}
\setlength{\belowdisplayskip}{4pt}
\setlength{\abovedisplayshortskip}{3pt}
\setlength{\belowdisplayshortskip}{3pt}
\renewcommand{\arraystretch}{0.95}

\newcommand{\pmv}[2]{#1{\scriptsize$\pm$#2}}

\begin{document}

\title{DSGD on RIPPER/FOIL-Induced Rules}

\author{
\IEEEauthorblockN{Sargis Vardanyan}
\IEEEauthorblockA{Institute for Informatics and\\
Automation Problems\\
American University of Armenia\\
Yerevan, Armenia\\
sargis\_vardanyan@edu.aua.am}
\and
\IEEEauthorblockN{Aik Tarkhanyan}
\IEEEauthorblockA{Department of Statistics\\
LMU Munich\\
Munich, Germany\\
aik.tarkhanyan@campus.lmu.de}
\and
\IEEEauthorblockN{Ashot Harutyunyan}
\IEEEauthorblockA{Yerevan State University\\
Institute for Informatics and\\
Automation Problems\\
Yerevan, Armenia\\
harutyunyan.ashot@ysu.am}
}

\maketitle

\begin{abstract}
Prior Dempster-Shafer gradient-descent classifiers primarily study evidential mass learning and prediction from a supplied rule representation. This paper evaluates DSGD-Auto under a cached induced-rule-pool protocol: FOIL-style and RIPPER-style inducers first generate rule-pool exports, and DSGD then learns evidential masses without adding, removing, or rewriting rules inside each export. We compare learned evidential aggregation against deterministic rule-only baselines and a Random Forest baseline on nine tabular datasets. Across this benchmark, DSGD-Auto improves several rule-based aggregates while preserving fired-rule inspection and exposing residual evidential uncertainty, although the gains vary across datasets and metrics.
\end{abstract}

\begin{IEEEkeywords}
evidential classification, Dempster-Shafer theory, rule induction, FOIL, RIPPER, tabular classification
\end{IEEEkeywords}

\noindent\textbf{Code:} \url{https://github.com/SargisVardanian/DST}

\section{Introduction}

Rule-based classification represents prediction with explicit if-then conditions rather than opaque scores. FOIL and RIPPER are supervised rule-induction methods that learn such rules from labeled data \cite{quinlan1990,cohen1995}. FOIL grows class-directed rules by adding informative literals, while RIPPER learns ordered rules and prunes them to improve generalization.

Dempster-Shafer theory provides a way to attach not only class support, but also residual uncertainty to rule-based evidence \cite{dempster1968,shafer1976}. In the setting used here, each rule receives a mass vector over the $C$ class singletons and the full frame $\Omega$. Following DSGD/DSGD++ \cite{penafiel2020,tarkhanyan2025}, the model learns per-rule mass parameters, combines the masses of fired rules, and converts the fused mass into class probabilities through the pignistic transform \cite{smets1994}. In the reported learned-Dempster rows, the residual mass assigned to $\Omega$ is redistributed uniformly over classes, i.e., $m^\star(\Omega)/C$.

The evaluation uses cached rule-pool exports. Rule induction is performed before mass learning, and each downstream aggregation method is evaluated on the rules available in a given export. Thus, the DSGD stage updates only evidential masses within that export. The FOIL-style and RIPPER-style inducers serve as study-specific fixed rule generators for the aggregation study, rather than as the main objects of comparison. Because some baseline conventions differ on no-fire samples, the study should be read as a practical benchmark of rule-use strategies rather than as a perfectly isolated aggregation-only ablation.

\section{Experimental Setup}

This evaluation focuses on supervised FOIL-style and RIPPER-style induced rule sets. For each dataset, the data are partitioned into train and test subsets according to the benchmark pipeline. Rule induction is performed only on the training subset, and the resulting rule set is cached before downstream aggregation is evaluated. DSGD trains only evidential mass parameters on top of the cached rules; it does not add, remove, or rewrite rule antecedents. Random Forest is evaluated under the same train/test partitioning protocol as an external non-rule baseline \cite{breiman2001}.

The present snapshot does not perform cross-validation or repeated train/test resampling. Consequently, the results should be interpreted as a benchmark comparison under the current partitioning protocol, not as split-robust evidence of superiority.

Reported values use the form mean $\pm$ standard deviation over the exported benchmark repetitions.

The tabular benchmark contains nine public classification datasets: adult, bank-full, BrainTumor, breast-cancer-wisconsin, df\_wine, dry-bean, gas\_drift, german, and magic-gamma. Standard preprocessing is applied before rule induction, including missing-value filtering and categorical encoding required by the rule inducers. The exact implementation details, rule-generator settings, and cached export generation logic are part of the released benchmark code. Since the present paper does not fully enumerate all inducer hyperparameters, conclusions about FOIL-versus-RIPPER induction quality should be treated cautiously.

The benchmark compares four ways of using the cached induced rules. First, an ordered-rule baseline uses only the first activated rule. Second, a first-hit probability wrapper converts that ordered prediction into a fixed artificial probability vector. Third, weighted vote combines all activated rules using support- and confidence-based rule weights. Fourth, learned Dempster fusion combines learned evidential mass vectors attached to the activated rules. These definitions are given explicitly below because the methods differ not only in aggregation, but also in their behavior when no rule is active.

Let $C$ be the number of classes. For a test sample $x$, let
\[
  \mathcal{A}(x)=(r_1,\ldots,r_m)
\]
be the ordered list of rules whose antecedents are satisfied by $x$. A rule is active if all literals in its antecedent are true for the sample. Each rule $r$ has an attached consequent class $\hat{y}(r)$. The order of $\mathcal{A}(x)$ is inherited from the emitted rule-pool export. Let $n_k$ be the number of training samples from class $k$, and define the normalized inverse-frequency class-balance vector
\[
  \pi_k=\frac{1/n_k}{\sum_{j=1}^{C}1/n_j}.
\]
This vector is used only as the no-fire probability fallback for weighted vote.

The ordered-rule baseline predicts from the first activated rule:
\[
  \hat{y}_{\mathrm{native}}(x)=
  \begin{cases}
    \hat{y}(r_1), & m>0,\\
    y_{\mathrm{default}}, & m=0.
  \end{cases}
\]
Here $y_{\mathrm{default}}$ is the majority class in the training subset. This baseline is mainly a label-producing rule system. Therefore, its exported NLL and ECE values should be interpreted as diagnostic wrapper values rather than as calibrated probabilities produced by the rule system itself.

The first-hit fixed-confidence baseline keeps the same hard prediction as the ordered-rule baseline, but assigns a fixed probability vector around it:
\[
  p^{\mathrm{fh}}_k(x)=
  \begin{cases}
    0.9, & k=\hat{y}_{\mathrm{native}}(x),\\
    \frac{0.1}{C-1}, & k\neq\hat{y}_{\mathrm{native}}(x).
  \end{cases}
\]
This wrapper is not learned and is not estimated from rule support counts. It is included only so that label-only ordered predictions can be evaluated with probability-based diagnostics.

The weighted-vote baseline uses all active rules. For a rule $r$, let $s_r$ be the number of training samples satisfying the antecedent of $r$ and belonging to $\hat{y}(r)$. Let $n_r$ be the number of training samples satisfying the same antecedent but belonging to a different class. The rule confidence proxy and vote weight are
\[
  q_r=\frac{s_r+1}{s_r+n_r+C},
  \qquad
  w_r=q_r\log(1+s_r).
\]
For a sample $x$, the class vote score is
\[
  v_k(x)=\sum_{r\in\mathcal{A}(x)}
  w_r\mathbf{1}\{\hat{y}(r)=k\}.
\]
If at least one rule is active, weighted vote normalizes these scores:
\[
  p^{\mathrm{wv}}_k(x)=\frac{v_k(x)}{\sum_{j=1}^{C}v_j(x)}.
\]
If no rule is active, it uses the class-balance fallback:
\[
  p^{\mathrm{wv}}_k(x)=\pi_k.
\]
The hard prediction is always obtained from the same probability vector:
\[
  \hat{y}_{\mathrm{wvote}}(x)=\arg\max_k p^{\mathrm{wv}}_k(x).
\]
Ties are broken by the fixed encoded class order. This fallback is a benchmark convention; it is not a Dempster-Shafer vacuous-mass state.

For learned Dempster fusion, each active rule carries a learned sparse mass vector. In the implementation used here, rule $r$ can assign positive singleton mass only to its attached consequent class $\hat{y}(r)$ and positive residual mass to $\Omega$; all other singleton masses are fixed to zero. We still write the vector in $C{+}1$ form,
\[
  m_r=[m_r(\{1\}),\dots,m_r(\{C\}),m_r(\Omega)].
\]
Its nonzero entries are restricted to $m_r(\{\hat{y}(r)\})$ and $m_r(\Omega)$. For two mass functions $m_1,m_2$, Dempster's normalized combination is
\[
  (m_1\oplus m_2)(A)=
  \frac{\sum_{B\cap D=A}m_1(B)m_2(D)}{1-\kappa},
  \quad A\neq\emptyset,
\]
where
\[
  \kappa=\sum_{B\cap D=\emptyset}m_1(B)m_2(D)
\]
is the conflict mass. Repeated fusion over the active rules gives
\[
  m^\star(x)=\bigoplus_{r\in\mathcal{A}(x)}m_r.
\]
The reported class probabilities are obtained by the pignistic readout:
\[
  p^{\mathrm{ds}}_k(x)=m^\star(\{k\})+\frac{m^\star(\Omega)}{C},
  \qquad
  \hat{y}_{\mathrm{dsgd}}(x)=\arg\max_k p^{\mathrm{ds}}_k(x).
\]
If no rule is active, learned Dempster fusion uses the vacuous mass vector $[0,\dots,0,1]$, which gives $p^{\mathrm{ds}}_k(x)=1/C$ and fused uncertainty $1$. Therefore, weighted vote and learned Dempster use different no-fire probability fallbacks. If no-fire samples are frequent, this difference can affect Accuracy, Macro-F1, NLL, and ECE.

To keep near-total-conflict fusion numerically well-defined, the implementation clips $\kappa$ to at most $1-10^{-9}$, lower-bounds $1-\kappa$ by $10^{-9}$, and renormalizes the resulting mass vector. The activation frequency of these safeguards is not separately analyzed in this snapshot.

\section{Validation and Discussion}

\subsection{Predictive Results}

The central question is whether learning evidential masses over an induced rule pool improves prediction relative to simpler ways of using the same cached rules. The comparison is not intended to show that FOIL or RIPPER induction is universally superior to other tabular learners. Because the methods use different no-fire conventions, part of the measured difference may also reflect fallback behavior on samples for which no rule fires.

Table~\ref{tab:dataset_snapshot} reports the best rule-based row on each dataset across the evaluated FOIL-style and RIPPER-style pipelines. This is a mixed-pipeline snapshot: it identifies which full rule-generation-plus-aggregation row performed best, but it should not by itself be read as an isolated aggregation-only comparison. Table~\ref{tab:avg_results} reports cross-dataset averages for each method family. For each dataset and method, metrics are first averaged over available exported repetitions; Table~\ref{tab:avg_results} then reports the unweighted mean of those dataset-level means across the nine datasets, with the displayed standard deviation computed over the corresponding repetition-level summaries exported by the benchmark pipeline. The uncertainty columns use the same averaging order. Mean uncertainty denotes average rule-level residual mass assigned to $\Omega$ over active-rule evaluations, and fused uncertainty denotes average residual mass $m^\star(\Omega)$ after rule fusion.

\begin{table}[!b]
  \caption{Best mixed-pipeline rule-based rows by dataset.}
  \label{tab:dataset_snapshot}
  \centering
  \footnotesize
  \setlength{\tabcolsep}{3.3pt}
  \renewcommand{\arraystretch}{1.00}
  \begin{tabular}{lcc}
    \toprule
    Dataset & Best Acc. & Best Macro-F1 \\
    \midrule
    adult & RIPPER + DSGD & RIPPER + DSGD \\
    bank-full & RIPPER + DSGD & RIPPER + DSGD \\
    BrainTumor & FOIL + wvote & FOIL + wvote \\
    breast-cancer-wisconsin & FOIL ordered & FOIL ordered \\
    df\_wine & FOIL + wvote & FOIL + wvote \\
    dry-bean & RIPPER + DSGD & RIPPER + DSGD \\
    gas\_drift & RIPPER + DSGD & RIPPER + DSGD \\
    german & RIPPER + DSGD & RIPPER + DSGD \\
    magic-gamma & RIPPER + DSGD & RIPPER + DSGD \\
    \bottomrule
  \end{tabular}
\end{table}

\begin{table}[!b]
  \caption{Average benchmark results. Values are reported as mean $\pm$ standard deviation and rounded to three decimals.}
  \label{tab:avg_results}
  \centering
  \scriptsize
  \setlength{\tabcolsep}{1.8pt}
  \renewcommand{\arraystretch}{1.05}
  \resizebox{\columnwidth}{!}{%
    \begin{tabular}{lcccccc}
      \toprule
      Method & Acc. & Macro-F1 & NLL & ECE & Mean unc. & Fused unc. \\
      \midrule
      \textbf{Random Forest}
        & \textbf{\pmv{.925}{.001}} & \textbf{\pmv{.890}{.002}}
        & \textbf{\pmv{.197}{.004}} & \textbf{\pmv{.024}{.002}}
        & --- & --- \\
      RIPPER + DSGD
        & \pmv{.910}{.002} & \pmv{.874}{.003}
        & \pmv{.291}{.007} & \pmv{.091}{.008}
        & \pmv{.620}{.010} & \pmv{.290}{.014} \\
      RIPPER + wvote
        & \pmv{.906}{.001} & \pmv{.860}{.001}
        & \pmv{1.060}{.088} & \pmv{.054}{.001}
        & --- & --- \\
      RIPPER ordered
        & \pmv{.841}{.029} & \pmv{.765}{.026}
        & \pmv{4.397}{.798} & \pmv{.159}{.029}
        & --- & --- \\
      RIPPER first-hit
        & \pmv{.841}{.029} & \pmv{.765}{.026}
        & \pmv{.572}{.084} & \pmv{.122}{.029}
        & --- & --- \\
      FOIL + DSGD
        & \pmv{.896}{.001} & \pmv{.847}{.002}
        & \pmv{.306}{.004} & \pmv{.087}{.004}
        & \pmv{.583}{.005} & \pmv{.329}{.007} \\
      FOIL + wvote
        & \pmv{.893}{.000} & \pmv{.837}{.000}
        & \pmv{1.416}{.000} & \pmv{.062}{.000}
        & --- & --- \\
      FOIL ordered
        & \pmv{.889}{.000} & \pmv{.829}{.000}
        & \pmv{3.080}{.001} & \pmv{.112}{.000}
        & --- & --- \\
      FOIL first-hit
        & \pmv{.889}{.000} & \pmv{.829}{.000}
        & \pmv{.410}{.000} & \pmv{.075}{.000}
        & --- & --- \\
      \bottomrule
    \end{tabular}%
  }
\end{table}

The per-dataset mixed-pipeline snapshot shows a non-uniform pattern. RIPPER + learned Dempster is the strongest rule-based row on six of the nine datasets for both Accuracy and Macro-F1. Weighted vote remains strongest on BrainTumor and df\_wine, while ordered FOIL rules remain strongest on breast-cancer-wisconsin. Because Table~\ref{tab:dataset_snapshot} compares complete FOIL-style and RIPPER-style pipelines, these wins combine rule-pool effects and aggregation effects.

Within each inducer family in Table~\ref{tab:avg_results}, RIPPER + learned Dempster improves over RIPPER weighted vote in mean Accuracy and Macro-F1, and FOIL + learned Dempster improves over FOIL weighted vote on both metrics. Random Forest gives the best overall predictive scores. DSGD-Auto is therefore best viewed as an interpretable evidential aggregation method rather than a replacement for Random Forest when only predictive performance matters.

Calibration is mixed. Learned Dempster fusion obtains substantially lower NLL than weighted vote for both FOIL and RIPPER, but weighted vote obtains lower aggregate ECE in this snapshot. Thus, learned Dempster fusion improves likelihood fit relative to weighted vote in these averages, but it does not improve calibration uniformly. A more precise analysis would require stratification by no-fire cases, active-rule counts, and high-conflict cases.

\begin{table}[!b]
  \caption{Average rule and aggregated-explanation lengths.}
  \label{tab:rule_lengths}
  \centering
  \footnotesize
  \setlength{\tabcolsep}{2.6pt}
  \renewcommand{\arraystretch}{1.00}
  \resizebox{\columnwidth}{!}{%
    \begin{tabular}{lcccc}
      \toprule
      Dataset & RIPPER & FOIL & RIPPER DS agg. & FOIL DS agg. \\
      \midrule
      adult & 5.01$\pm$0.18 & 5.73$\pm$0.00 & 5.97$\pm$0.43 & 6.15$\pm$0.00 \\
      bank-full & 4.81$\pm$0.00 & 5.43$\pm$0.00 & 8.50$\pm$0.00 & 9.83$\pm$0.00 \\
      BrainTumor & 2.53$\pm$0.00 & 3.80$\pm$0.00 & 5.28$\pm$0.00 & 6.27$\pm$0.00 \\
      breast-cancer-wisconsin & 1.31$\pm$0.00 & 1.96$\pm$0.00 & 3.69$\pm$0.00 & 4.19$\pm$0.01 \\
      df\_wine & 1.58$\pm$0.20 & 3.88$\pm$0.00 & 1.69$\pm$0.15 & 6.92$\pm$0.00 \\
      dry-bean & 5.29$\pm$0.00 & 6.22$\pm$0.00 & 10.84$\pm$0.00 & 11.96$\pm$0.00 \\
      gas\_drift & 3.98$\pm$0.00 & 4.14$\pm$0.00 & 8.44$\pm$0.00 & 8.55$\pm$0.00 \\
      german & 3.00$\pm$0.00 & 4.04$\pm$0.00 & 4.74$\pm$0.02 & 4.16$\pm$0.03 \\
      magic-gamma & 5.48$\pm$0.00 & 6.54$\pm$0.00 & 7.99$\pm$0.00 & 8.09$\pm$0.00 \\
      \bottomrule
    \end{tabular}%
  }
\end{table}

Table~\ref{tab:rule_lengths} reports the average length of the induced rules and the average length of the Dempster-Shafer aggregated explanations on the test subset. The RIPPER and FOIL columns measure the average number of antecedent literals in the original induced rules. The DS aggregated columns measure the average number of literals retained in the post-hoc explanation produced after the active rules are fused and their positive evidence is summarized.

The aggregated explanations are generally longer than individual induced rules because a Dempster-Shafer prediction may depend on several active rules rather than on a single first-hit rule. The aggregated explanation should therefore be interpreted as a compact post-hoc summary, not as a faithful algebraic decomposition of Dempster's rule.

\subsection{Worked Example}

The following example illustrates an aggregated explanation for one inspected sample from the Adult dataset. In this binary example, coordinate~1 corresponds to \code{<=50K}, and coordinate~2 corresponds to \code{>50K}. Five rules are active:

\begin{center}
  \scriptsize
  \setlength{\tabcolsep}{2.2pt}
  \begin{tabular}{@{}l p{0.52\columnwidth} l@{}}
    \toprule
    Rule & Antecedent excerpt & $m_r$ \\
    \midrule
    $R_{84}$ & age $>53$, capital gain $<5013$, education $=$ HS-grad, $\ldots$ & $[.772,.000,.228]$ \\
    $R_{199}$ & age $>30$, fnlwgt $>1.304{\times}10^{5}$, occupation $=$ Prof-specialty, $\ldots$ & $[.000,.447,.553]$ \\
    $R_{164}$ & age $>31$, hours per week $>32$, occupation $=$ Prof-specialty, $\ldots$ & $[.000,.307,.693]$ \\
    $R_{119}$ & age $>57$, capital gain $<5013$, hours per week $<42$, $\ldots$ & $[.299,.000,.701]$ \\
    $R_{169}$ & marital status $=$ Married-civ-spouse & $[.197,.000,.803]$ \\
    \bottomrule
  \end{tabular}
\end{center}

Each row shows a rule antecedent and the learned mass vector assigned to that rule. For a binary task, the first two coordinates represent support for the two class singletons, and the last coordinate represents residual uncertainty $m_r(\Omega)$. Under the sparse parameterization used here, a rule assigns singleton mass only to its attached consequent class.

The rule masses are fused for prediction by Dempster's rule. The aggregated explanation below is post-hoc: it summarizes which input conditions are most associated with the final prediction under the chosen explanation heuristic, but it does not change the prediction itself. Let $\hat{y}\equiv\hat{y}_{\mathrm{dsgd}}(x)$ be the final predicted class, and let $y_r\equiv\hat{y}(r)$ be the class label attached to active rule $r$. The explanation uses the normalized active-rule weight
\[
  \rho_r=\frac{w_r}{\sum_{r'\in\mathcal{A}(x)}w_{r'}},
\]
where $w_r$ is the support- and confidence-based rule weight used in weighted vote. If all active-rule weights are zero, $\rho_r$ is set uniformly over active rules. The explanation-only contribution of rule $r$ is
\[
  \alpha_r=\rho_r
  \begin{cases}
    m_r(\{\hat{y}\}), & y_r=\hat{y},\\[2pt]
    -\tfrac{1}{2}m_r(\{y_r\}), & y_r\neq\hat{y}.
  \end{cases}
\]
This linear form avoids an unintended squaring of rule commitment under the sparse parameterization. The negative half-weight is a heuristic penalty used only for explanation scoring; it is not part of Dempster fusion. Only positive $\alpha_r$ values are propagated to literals. If rule $r$ has $L_r$ antecedent literals, each receives $\alpha_r/L_r$. Repeated literals are then merged; equality conditions retain the most-supported value, while numeric bounds retain the tightest interval implied by contributing rules.

For the inspected sample, the aggregated explanation is:
\begin{equation*}
\begin{gathered}
\text{aggregated explanation}\approx \\
.150[\text{age} > 57] \\
{}+ .150[32 < \text{hours/week} < 40] \\
{}+ .145[\text{marital} = \text{Married-civ-spouse}] \\
{}+ .130[130420 < \text{fnlwgt} < 172960] \\
{}+ .116[\text{capital gain} < 5013] + \cdots .
\end{gathered}
\end{equation*}
The notation is intentionally compact: each bracketed condition is a literal or a merged literal range, and the coefficient is its positive post-hoc explanation score after active-rule contributions have been distributed across literals.

The final fused output is:
\[
  p^{\mathrm{ds}}=[.776,.224],
  \qquad
  m^\star=[.723,.171,.106].
\]
The pignistic output assigns probability $.776$ to \code{<=50K} and $.224$ to \code{>50K}; the fused mass is $[.723,.171,.106]$. The prediction is therefore dominated by \code{<=50K} evidence, with limited opposing support and residual uncertainty. The explanation summarizes positively scored fired-rule conditions, but it is not a causal proof.

\subsection{Limitations}

The main limitation is that Dempster-Shafer aggregation does not create new rule structure. DSGD-Auto can only reweight and combine the evidence supplied by the induced rule pool. If the rule pool misses important regions of the input space, contains weak antecedents, or encodes misleading consequents, learned masses can reduce but not remove this structural limitation.

The second limitation is the independence assumption implicit in Dempster-style fusion. Fired rules are often not independent: they may share literals, cover overlapping samples, or express similar decision patterns. In such cases, combining them as separate pieces of evidence can overemphasize repeated information. The learned uncertainty mass partly softens this effect, but it does not fully solve dependence between rules.

A further limitation is conflict handling. Dempster's normalized rule redistributes conflict among non-empty hypotheses. This can be useful when several partially reliable rules disagree, but it can also amplify confident evidence when the active rules are strongly contradictory. Therefore, high-conflict cases require separate inspection rather than being treated as ordinary aggregated predictions.

No-fire behavior is also a theoretical boundary of the model. When no rule fires, learned Dempster fusion returns a vacuous mass and therefore a uniform pignistic distribution. This is a conservative uncertainty state, not a learned extrapolation outside the covered rule regions. Thus, DSGD-Auto is naturally strongest when the induced rules provide meaningful coverage of the test space.

Finally, the method should be understood as evidential aggregation over a fixed symbolic representation, not as a full theory of rule induction. The FOIL-style and RIPPER-style generators determine which symbolic evidence is available; DSGD-Auto only learns how much support and residual uncertainty to assign to that evidence. Future work should therefore study rule dependence, conflict behavior, coverage, and uncertainty calibration as theoretical properties of the combined rule-and-evidence system.

\section{Conclusion}

This work evaluates DSGD-Auto as learned evidential aggregation over cached FOIL/RIPPER-style rule pools. The method does not change the symbolic rule structure; instead, it learns how much class support and residual uncertainty each fired rule should contribute.

The main result is that learned Dempster-Shafer fusion can improve rule-based prediction while preserving inspectability of fired rules, rule masses, fused uncertainty, and post-hoc explanations. Its value is therefore not only predictive, but also structural: it connects symbolic rule evidence with uncertainty-aware aggregation.

The method is not a universal replacement for strong tabular ensembles. Its behavior depends on the quality, coverage, dependence, and conflict structure of the induced rule pool. Future work should study these theoretical properties more directly, especially rule dependence, conflict normalization, no-fire uncertainty, and calibration of learned evidential masses.

\begin{thebibliography}{99}

\bibitem{dempster1968}
A. P. Dempster, ``A generalization of Bayesian inference,'' \emph{Journal of the Royal Statistical Society: Series B}, vol. 30, no. 2, pp. 205--232, 1968.

\bibitem{shafer1976}
G. Shafer, \emph{A Mathematical Theory of Evidence}. Princeton University Press, 1976.

\bibitem{smets1994}
P. Smets and R. Kennes, ``The transferable belief model,'' \emph{Artificial Intelligence}, vol. 66, no. 2, pp. 191--234, 1994.

\bibitem{quinlan1990}
J. R. Quinlan, ``Learning logical definitions from relations,'' \emph{Machine Learning}, vol. 5, pp. 239--266, 1990.

\bibitem{cohen1995}
W. W. Cohen, ``Fast effective rule induction,'' in \emph{Machine Learning Proceedings 1995}, pp. 115--123, Morgan Kaufmann, 1995.

\bibitem{penafiel2020}
S. Pe\~nafiel, N. Baloian, H. Sanson, and J. A. Pino, ``Applying Dempster-Shafer theory for developing a flexible, accurate and interpretable classifier,'' \emph{Expert Systems with Applications}, vol. 157, p. 113262, 2020.

\bibitem{tarkhanyan2025}
A. Tarkhanyan and A. Harutyunyan, ``DSGD++: Performance and robustness improvements for the Dempster-Shafer Gradient Descent classifier,'' arXiv:2507.00453, 2025.

\bibitem{breiman2001}
L. Breiman, ``Random forests,'' \emph{Machine Learning}, vol. 45, no. 1, pp. 5--32, 2001.

\end{thebibliography}

\end{document}
